% ============================================================ antisymmetry and spin
\section{Spin \& antisymmetry}

\begin{frame}{Swapping two electrons}{why the exchange eigenvalue is $\pm1$}
  \begin{columns}[c]
    \column{0.62\textwidth}
    Define $\hat P_{12}\Psi(x_1,x_2)=\Psi(x_2,x_1)$, \quad $x=(\mathbf r,\sigma)$.
    {\begin{align*}
      \stepno{1}\ &\hat P_{12}^2=\hat 1 &&\why{swapping twice restores the order}\\
      \uncover<2->{\stepno{2}\ }&\uncover<2->{\hat P_{12}\Psi=\lambda\Psi\ \Rightarrow\ \Psi=\lambda^2\Psi} &&\uncover<2->{\why{apply $\hat P_{12}$ twice}}\\
      \uncover<3->{\stepno{3}\ }&\uncover<3->{\lambda=\pm1} &&\uncover<3->{\why{since $\lambda^2=1$}}\\
      \uncover<4->{\stepno{4}\ }&\uncover<4->{[\Hop,\hat P_{12}]=0} &&\uncover<4->{\why{$\Hop$ treats all electrons alike}}
    \end{align*}}
    \column{0.34\textwidth}
    \centering
    \begin{tikzpicture}[scale=1.0]
      \fill[noircyan] (0,0) circle (0.12) node[below=3pt,font=\scriptsize]{1};
      \fill[noirrose] (2,0) circle (0.12) node[below=3pt,font=\scriptsize]{2};
      \draw[-{Stealth[length=5pt]},noirgold,line width=0.8pt] (0.2,0.15) to[bend left=50] (1.8,0.15);
      \draw[-{Stealth[length=5pt]},noirgold,line width=0.8pt] (1.8,-0.15) to[bend left=50] (0.2,-0.15);
      \node[font=\small,text=noirgold] at (1,0.95) {$\hat P_{12}$};
      \onslide<5->{\node[font=\large,text=noirrose] at (1,-1.15) {$\lambda=-1$};}
    \end{tikzpicture}
  \end{columns}
  \onslide<5->{\small\cgold{so eigenstates have definite $\lambda$; for electrons nature picks $-1$.} The paper: ``the wavefunction for a
    system of fermions must be antisymmetric with respect to the interchange of any two particles'' (p.\ 441)\par}
  \symbolstrip{$x_i=(\mathbf r_i,\sigma_i)$: all coordinates of electron $i$, position $\mathbf r_i$ plus spin coordinate $\sigma_i$ ($\uparrow$ or $\downarrow$)
    \quad $\hat P_{12}$: operator exchanging electrons 1 and 2 \quad $\lambda$: its eigenvalue \quad $[\hat A,\hat B]=\hat A\hat B-\hat B\hat A$
    \quad fermions: half-integer-spin particles (electrons)}
\end{frame}

\begin{frame}{Slater determinants}{antisymmetry built in}
  \[
    D[\chi_1,\chi_2]=\begin{vmatrix}\chi_1(1)&\chi_2(1)\\\chi_1(2)&\chi_2(2)\end{vmatrix}
    =\chi_1(1)\chi_2(2)-\chi_2(1)\chi_1(2)
  \]
  \vskip-0.2cm
  \explain{rows label electrons, columns label spin orbitals; the determinant mixes every assignment with a sign}
  \begin{itemize}
    \item<2-> swap electrons $\Leftrightarrow$ swap two \emph{rows} $\Rightarrow$ $D\to-D$ \hfill\cgreen{antisymmetric}
    \item<3-> two equal spin orbitals $\Leftrightarrow$ two equal \emph{columns} $\Rightarrow$ $D=0$ \hfill\cgold{Pauli exclusion}
    \item<4-> normalised: $\ket{\chi_1\chi_2}=\tfrac1{\sqrt2}D[\chi_1,\chi_2]$
  \end{itemize}
  \onslide<5->{
  \begin{keyresult}
    Pauli's 1925 rule (no two electrons with all four quantum numbers equal; p.\ 448) is a \emph{consequence}
    of antisymmetry inside the orbital model, not a separate law.
  \end{keyresult}}
  \symbolstrip{$\chi(x)$: spin orbital $=$ spatial orbital $\times$ spin function, e.g.\ $\chi=\pa(\mathbf r)\,\alpha(\sigma)$ \quad
    $\chi_1(2)$: spin orbital 1 at electron 2's coordinates \quad $\tfrac1{\sqrt2}$: $D$ is two orthogonal products of norm 1, so $\braket DD=2$}
\end{frame}

\begin{frame}{Spin operators for two electrons}
  \begin{columns}[T]
    \column{0.47\textwidth}
    \cgold{One electron} ($s=\tfrac12$):
    {\begin{align*}
      \hat s_z\alpha&=+\tfrac12\alpha, & \hat s_z\beta&=-\tfrac12\beta\\
      \hat s_+\beta&=\alpha, & \hat s_+\alpha&=0\\
      \hat s_-\alpha&=\beta, & \hat s_-\beta&=0
    \end{align*}}
    \onslide<2->{\cgold{Two electrons:} add each electron's operator\par
      $\hat S_z=\hat s_{z}(1)+\hat s_{z}(2)$,\ \ $\hat S_\pm=\hat s_\pm(1)+\hat s_\pm(2)$}
    \column{0.5\textwidth}
    \onslide<3->{\cgold{A useful identity:}}
    {\begin{align*}
      \uncover<3->{\hat S_-\hat S_+}&\uncover<3->{=(\hat S_x-i\hat S_y)(\hat S_x+i\hat S_y)} &&\uncover<3->{\why{definition of $\hat S_\pm$}}\\
      &\uncover<4->{=\hat S_x^2+\hat S_y^2+i[\hat S_x,\hat S_y]} &&\uncover<4->{\why{multiply out}}\\
      &\uncover<5->{=\hat S^2-\hat S_z^2+i(i\hat S_z)} &&\uncover<5->{\why{$[\hat S_x,\hat S_y]=i\hat S_z$}}\\
      &\uncover<6->{=\hat S^2-\hat S_z^2-\hat S_z} &&\uncover<6->{\why{$i\cdot i=-1$}}
    \end{align*}}
    \onslide<7->{\[
      \boxed{\hat S^2=\hat S_-\hat S_++\hat S_z^2+\hat S_z}
    \]}
  \end{columns}
  \symbolstrip{$\alpha,\beta$: spin up / down of one electron \quad $\hat s_+$ flips $\downarrow\to\uparrow$, $\hat s_-$ flips $\uparrow\to\downarrow$ (0 if impossible)
    \quad $\hat S_x,\hat S_y,\hat S_z,\hat S^2$: total spin components and squared magnitude \quad $\hbar=1$}
\end{frame}

\begin{frame}{$\hat S^2$ on the four spin products (1/2)}
  \begin{columns}[T]
    \column{0.48\textwidth}
    \cgold{$\alpha_1\alpha_2$} \ \cmuted{(both up)}
    {\begin{align*}
      \hat S_+\alpha_1\alpha_2&=0 &&\why{nothing left to raise}\\
      \uncover<2->{\hat S_z\alpha_1\alpha_2}&\uncover<2->{=(\tfrac12+\tfrac12)\alpha_1\alpha_2} &&\uncover<2->{\why{each adds $+\tfrac12$}}\\
      \uncover<3->{\hat S^2\alpha_1\alpha_2}&\uncover<3->{=(0+1+1)\,\alpha_1\alpha_2} &&\uncover<3->{\why{$\hat S_-\hat S_+ +\hat S_z^2+\hat S_z$}}\\
      &\uncover<4->{=\cgreen{2}\,\alpha_1\alpha_2}
    \end{align*}}
    \onslide<4->{\small $S(S+1)=2\Rightarrow S=1,\ M_S=+1$}
    \column{0.48\textwidth}
    \onslide<5->{\cgold{$\beta_1\beta_2$} \ \cmuted{(both down)}}
    {\begin{align*}
      \uncover<5->{\hat S_+\beta_1\beta_2}&\uncover<5->{=\alpha_1\beta_2+\beta_1\alpha_2} &&\uncover<5->{\why{raise 1 or 2}}\\
      \uncover<6->{\hat S_-(\ldots)}&\uncover<6->{=\beta_1\beta_2+\beta_1\beta_2} &&\uncover<6->{\why{lower the raised one}}\\
      \uncover<7->{(\hat S_z^2+\hat S_z)\beta_1\beta_2}&\uncover<7->{=(1-1)\beta_1\beta_2} &&\uncover<7->{\why{$M_S=-1$}}\\
      \uncover<8->{\hat S^2\beta_1\beta_2}&\uncover<8->{=\cgreen{2}\,\beta_1\beta_2}
    \end{align*}}
    \onslide<8->{\small $S=1,\ M_S=-1$}
  \end{columns}
  \symbolstrip{$\alpha_1\equiv\alpha(\sigma_1)$: electron 1 is spin up \quad $\alpha_1\beta_2$: electron 1 up, electron 2 down
    \quad eigenvalue of $\hat S^2$ is $S(S+1)$ \quad $M_S$: eigenvalue of $\hat S_z$ (sum of the two $\pm\tfrac12$)}
\end{frame}

\begin{frame}{$\hat S^2$ on the four spin products (2/2)}
  \cgold{$\alpha_1\beta_2$ and $\beta_1\alpha_2$} \quad\cmuted{($M_S=0$; neither is an eigenfunction on its own)}
  {\begin{align*}
    \hat S_+\alpha_1\beta_2&=\alpha_1\alpha_2 &
    \uncover<2->{\hat S_-\alpha_1\alpha_2}&\uncover<2->{=\beta_1\alpha_2+\alpha_1\beta_2}\\
    \uncover<3->{\hat S_z\alpha_1\beta_2}&\uncover<3->{=(\tfrac12-\tfrac12)\,\alpha_1\beta_2=0} &
    \uncover<4->{\Rightarrow\ \hat S^2\alpha_1\beta_2}&\uncover<4->{=\alpha_1\beta_2+\beta_1\alpha_2}\\
    &&\uncover<5->{\hat S^2\beta_1\alpha_2}&\uncover<5->{=\alpha_1\beta_2+\beta_1\alpha_2\quad\why{same steps}}
  \end{align*}}
  \explain<4->{$\hat S^2$ turns $\alpha_1\beta_2$ into a \emph{different} function: it is not an eigenfunction}
  \onslide<6->{Add and subtract the two results:}
  \[
    \uncover<6->{\hat S^2(\alpha_1\beta_2\pm\beta_1\alpha_2)=(1\pm1)\,(\alpha_1\beta_2\pm\beta_1\alpha_2)}
  \]
  \onslide<7->{\centering
  \begin{tikzpicture}
    \node[draw=noirrose,rounded corners,fill=noirpanel,inner sep=5pt] at (0,0) {$+$:\ eigenvalue $2$\ \ $\Rightarrow S=1$};
    \node[draw=noirblue,rounded corners,fill=noirpanel,inner sep=5pt] at (5.2,0) {$-$:\ eigenvalue $0$\ \ $\Rightarrow S=0$};
  \end{tikzpicture}\par}
  \symbolstrip{$\hat S_+$ acts on electron 1 or 2 in turn (it is a sum); only the $\beta$ of electron 2 can be raised here \quad
    $(1\pm1)$: $1+1=2$ for the sum, $1-1=0$ for the difference}
\end{frame}

\begin{frame}{Triplet and singlet spin functions}
  \begin{columns}[c]
    \column{0.6\textwidth}
    \begin{keyresult}[The four two-electron spin states]
      \vskip-0.6em
      {\begin{align*}
        \Theta_{1,+1}&=\alpha_1\alpha_2\\
        \Theta_{1,0}&=\tfrac1{\sqrt2}(\alpha_1\beta_2+\beta_1\alpha_2)\\
        \Theta_{1,-1}&=\beta_1\beta_2\\
        \Theta_{0,0}&=\tfrac1{\sqrt2}(\alpha_1\beta_2-\beta_1\alpha_2)
      \end{align*}}
    \end{keyresult}
    \explain{$\tfrac1{\sqrt2}$: the two products are orthonormal, so the sum has norm$^2$ 2}
    \column{0.36\textwidth}
    \begin{itemize}
      \item<2-> triplet spin functions are \cgold{symmetric} under $1\leftrightarrow2$
      \item<3-> the singlet is \cgold{antisymmetric}
      \item<4-> total must be antisymmetric, so
    \end{itemize}
    \onslide<4->{
    \begin{center}
      \small
      \begin{tabular}{lcc}
        \toprule & space & spin\\\midrule
        triplet & \crose{anti} & sym\\
        singlet & \cblue{sym} & anti\\\bottomrule
      \end{tabular}
    \end{center}}
  \end{columns}
  \onslide<5->{\centering\small\tick\ \cmuted{T1: $\hat S^2,\hat S_z$ as $4\times4$ matrices in SymPy; all eigen-relations exact.}\par}
  \symbolstrip{$\Theta_{S,M_S}$: two-electron spin function with total spin $S$ and projection $M_S$ \quad
    symmetric: unchanged when $1\leftrightarrow2$; antisymmetric: changes sign \quad (anti)$\times$(sym) $=$ anti, as required}
\end{frame}
